%% BioCompiler (VSTTE submission) — Sections 7–9
%% Implementation, Evaluation, Related Work
%%
%% Compile as part of the main document (requires commands defined in
%% main.tex: \PASS, \FAIL, \UNCERTAIN, \tvl, \tand, \seq, \pred,
%% \verdict, \judg, \GOLD, \SILVER, \BRONZE)

% =================================================================
\FloatBarrier
\section{Implementation}\label{sec:impl}

BioCompiler is implemented as a Python~3.11 package comprising
166,800~lines of source code, 8,495~lines of Lean~4 proof code, and
10,779~lines of documentation.  The implementation is organized into four
main subsystems: the type checker and predicate library, the integrated
optimizer, the analysis engine integrations, and the certificate
infrastructure.  We describe each in turn.

\subsection{Package Structure}\label{sec:package}

The Python package is organized into the following modules:

\begin{table}[!htbp]
\centering
\caption{BioCompiler Python package structure. LOC counts include
         inline comments but exclude blank lines.}
\label{tab:packageloc}
\begin{tabular}{lrp{4.2cm}}
\toprule
\textbf{Module} & \textbf{LOC} & \textbf{Role} \\
\midrule
\texttt{biocompiler.optimizer}        & 27,924 & Integrated constraint-solving optimizer \\
\texttt{biocompiler.mhc\_binding\_db}   & 21,959 & MHC binding database and scoring \\
\texttt{biocompiler.benchmarking}     & 16,491 & Benchmarking and comparison infrastructure \\
\texttt{biocompiler.solver}           & 15,115 & CSP solver (OR-Tools, greedy) \\
\texttt{biocompiler.organisms}        &  9,697 & Codon usage tables for 30 organisms (5 with extended constraint profiles) \\
\texttt{biocompiler.engines}          &  9,714 & FFI engine wrappers (AlphaFold, FoldX, etc.) \\
\texttt{biocompiler.immunogenicity}  &  8,783 & Immunogenicity prediction and deimmunization \\
\texttt{biocompiler.validation}       &  8,434 & Dataset and literature validation \\
\texttt{biocompiler.provenance}       &  7,064 & Certificate generation and provenance tracking \\
\texttt{biocompiler.type\_system}      &  7,043 & 43 type predicates and type checking \\
\texttt{biocompiler.structure}        &  5,262 & Protein structure parsing (PDB, AlphaFold) \\
\texttt{biocompiler.sequence}        &  5,157 & Scanner, splicing (NDFST), MaxEntScan \\
\texttt{biocompiler.export}           &  4,278 & GenBank, FASTA, SBOL2, JSON export \\
\texttt{biocompiler.expression}       &  3,635 & mRNA stability, tAI, CAI computation \\
\texttt{biocompiler.api}              &  3,445 & REST API and Jupyter integration \\
\texttt{biocompiler.biosecurity}      &  3,258 & Biosecurity screening (BLAST, keyword) \\
\texttt{biocompiler.infrastructure}  &  3,439 & Logging, caching, configuration \\
\texttt{biocompiler.cli}              &  2,040 & Command-line interface \\
\texttt{biocompiler.application}      &  1,398 & Application-level orchestration \\
\texttt{biocompiler.shared}           &  1,371 & Five-valued logic, types, constants \\
\midrule
\textbf{Total}                        & \textbf{~166,800} & \\
\bottomrule
\end{tabular}
\end{table}

The test suite comprises 13,359~tests totaling ~169,000~lines of test code,
including 500~property-based tests using the Hypothesis
framework~\cite{maciver2019hypothesis} (\S\ref{sec:pbt}).

\subsection{Integrated Constraint-Solving Optimizer}\label{sec:optimizer}

The optimizer is a single integrated constraint-solving routine
(\texttt{integrated\_optimize}) that combines a greedy forward pass
with multi-pass cleanup.  It is explicitly \emph{not} a single
integrated pass that satisfies all constraints simultaneously: the
forward pass places codons left-to-right and can only evaluate
backward-looking and within-codon constraints, because the next codon
has not been placed yet.  Right-boundary and residual cross-codon
violations (GT/CG dinucleotides, ATTTA motifs spanning codon
boundaries, T-runs, restriction sites) are resolved in a dedicated
multi-pass cleanup that has full visibility of both neighbours.  The
high-level progression is:

\begin{enumerate}[leftmargin=*,itemsep=1pt,topsep=2pt]
\item \textbf{Back-translate}: sort synonymous codons by CAI
      (best first) for the target organism's codon usage table.
\item \textbf{Greedy forward pass} ($n$ positions, left-to-right):
      at each position try the best-CAI codon, check backward-looking
      constraints against the previous codon plus within-codon
      constraints (GT, CG, ATTTA, T-run, restriction site, premature
      stop); if violations, try the next-best codon, selecting the
      highest-CAI codon that satisfies the backward and within-codon
      constraints.
\item \textbf{Multi-pass cleanup} (up to 3~iterations): scan every
      codon with full left and right neighbour visibility; fix
      GT/CG/ATTTA/T-run violations by swapping to the best-CAI
      alternative codon.
\item \textbf{CpG sweep} (aggressive mode): targeted swaps for any
      remaining CG dinucleotides, including boundary CGs between codons.
\item \textbf{GC adjustment}: if sequence GC is outside $[lo, hi]$,
      swap codons toward the target window while preserving constraints.
\item \textbf{Verify and certify}: evaluate all 43~type predicates on
      the final sequence and produce a graduated certificate
      (\S\ref{sec:certs}).
\end{enumerate}

The integrated optimizer is deterministic (canonical codon ordering)
and produces the same output given the same input.  After step~6, if
any predicate returns \FAIL{} due to GT-mandatory positions, the
type-directed mutagenesis loop (\S\ref{sec:mut}) is invoked.

\subsection{Optimization Strategy}\label{sec:strategies}

The integrated optimizer always performs a greedy highest-CAI forward
pass followed by multi-pass cleanup; there is no separate strategy
dispatch.  The forward pass has complexity $O(n \times k)$ where $n$
is the number of codons and $k \le 6$ is the number of synonymous
alternatives per amino acid; the cleanup is $O(n \times k \times m)$
with $m \le 3$ iterations.  This replaces the earlier sequential
greedy optimizer that required 12+~separate passes, reducing the
constant factor while preserving all constraint invariants.

When a constraint-satisfaction solver is available
(\texttt{engine\_ortools}, backed by OR-Tools), the optimizer may
optionally delegate small sequences to the CSP backend for a globally
optimal assignment; if the backend is unavailable or the problem is
infeasible within the time budget, the greedy path takes over
seamlessly.  The greedy forward pass is the default and only
guaranteed execution path.

\subsection{Four Analysis Engines}\label{sec:engines}

BioCompiler integrates four external analysis engines through the SLOT-fill
protocol (\S\ref{sec:slot}).  Each engine provides domain-specific
predictions that enrich the type system's verdicts.  Crucially, engine
output feeds into \emph{SLOT-annotated} predicates only; the soundness
proof (Theorem~\ref{thm:sound}) is independent of engine accuracy.

\paragraph{ESMFold (protein structure prediction).}
ESMFold~\cite{lin2023esm} provides per-residue pLDDT confidence scores
for the predicted 3D structure.  BioCompiler calls the ESMFold API with
exponential-backoff retry (3~attempts, 5-second initial delay).  When the
API is unavailable, a heuristic fallback combines Kyte--Doolittle
hydrophobicity profiling, Chou--Fasman secondary structure
prediction~\cite{chou1974prediction}, and charge distribution analysis
to produce a low-confidence structural estimate.
% NOTE (heuristic estimates, not validated against experimental data):
The fallback's per-residue pLDDT scores are capped at 55 (vs.~ESMFold's
0--100 range) and the Chou--Fasman component has an expected Q3 accuracy
of ~50--60\% on standard benchmarks.  No systematic correlation study
against ESMFold has been performed; the heuristic is intended only as a
last resort to avoid blocking the optimizer when the API is unreachable.
Results are flagged with \texttt{method=\"heuristic\_fallback\"} and
\texttt{confidence~<~0.5}.

\paragraph{FoldX (stability estimation).}
FoldX~\cite{schymkowitz2005foldx} estimates the folding free energy
$\Delta\Delta G$ of point mutations.  BioCompiler invokes FoldX via its
command-line interface with the \texttt{--command BuildModel} option.  The
CLI version reports $\Delta\Delta G$ with a stated accuracy of
$\pm 1$~kcal/mol; our empirical measurements on 2,847~mutations from the
ProTherm database~\cite{kumar2006protherm} yield a mean absolute error
(MAE) of 3.4~kcal/mol, consistent with prior independent
evaluations~\cite{burghardt2020foldx}.  When FoldX is unavailable, a
heuristic based on BLOSUM62 substitution scores provides a rough
stability proxy (MAE~4.1~kcal/mol).

\paragraph{CamSol (solubility and aggregation).}
CamSol~\cite{sormanni2015camsol} predicts intrinsic solubility profiles
and identifies aggregation-prone regions.  For intrinsically disordered
proteins (IDPs), BioCompiler augments CamSol's default hydrophobicity
scale (Wimley--White~\cite{wimley1996experimental}) with the Urry
hydrophobicity scale~\cite{urry1992hydrophobicity}, which better captures
the conformational flexibility of IDPs.
% NOTE (heuristic estimates, not validated against large experimental datasets):
On a curated benchmark of 21~proteins, the intrinsic-mode classifier
achieves ~85.7\% overall accuracy with ~66.7\% sensitivity for
aggregation-prone/IDP sequences (improved from ~33.3\% using
Wimley--White alone by adding the Urry scale for predicted IDPs).
These figures are from a small internal benchmark and have not been
validated on larger experimental datasets; they should be interpreted
as preliminary heuristic estimates.  The Urry scale is selected
automatically when the predictor detects IDP characteristics (low
hydrophobic cluster density, high net charge).

\paragraph{NetMHCpan (immunogenicity).}
NetMHCpan~\cite{reynisson2020netmhcpan} predicts peptide--MHC binding
affinity, the primary determinant of T-cell epitope immunogenicity.
BioCompiler queries the NetMHCpan REST API with a configurable panel of
MHC alleles (default: 12~common HLA class~I alleles).  API responses are
cached locally with a 24-hour TTL to avoid redundant queries across
benchmark runs.  The API achieves AUC values of 0.85--0.95 across MHC-I
alleles~\cite{reynisson2020netmhcpan}.  When the API is unavailable, a
position-specific scoring matrix (PSSM) fallback trained on the IEDB
database~\cite{vita2019iedb} provides predictions with AUC
0.60--0.75---significantly less accurate but sufficient to avoid
certifying sequences with strong immunogenicity signals.

\subsection{SLOT-Fill Protocol}\label{sec:slot}

The SLOT-fill protocol mediates all communication between BioCompiler and
external engines.  SLOTs are typed placeholders in the IR that are filled
by engine output at runtime.  The protocol supports three operating
modes:

\begin{description}[leftmargin=!,labelwidth=2.8cm,itemsep=1pt,topsep=1pt]
\item[SLOT \textsc{Skip}:]
The engine is not invoked; all SLOT-dependent predicates return
\UNCERTAIN.  This mode is used when engine availability is not
guaranteed (e.g., offline deployment).

\item[SLOT \textsc{Annotate}:]
The engine is invoked and its output annotates the IR, but
SLOT-dependent predicates still return \UNCERTAIN.  Annotations are
visible to the optimizer but not to the type checker, ensuring
certificate validity is engine-independent.

\item[SLOT \textsc{Verified}:]
The engine is invoked and SLOT-dependent predicates may return \PASS{}
or \FAIL{} based on engine output.  This mode trades engine-independence
for stronger verdicts.  Of the 43~predicates, 19~are SLOT
\textsc{Annotate}/\textsc{Skip} predicates whose verdicts are
independent of engine output; the remaining 17~core predicates plus
7~extended diagnostic predicates cover the full system;
19~SLOT \textsc{Verified}
predicates (e.g., \emph{StructureConserved}, \emph{NoImmunogenicEpitope})
that may issue \PASS{} verdicts contingent on engine accuracy.
\end{description}

The SLOT architecture ensures a clean separation: the soundness proof
covers all 36~predicates in the Lean4 model unconditionally, and
the 19~SLOT-\textsc{Verified} predicates are explicitly flagged in the
certificate with their engine dependency.

\subsection{Refinement Theorem}

The Lean~4 proof includes a refinement theorem relating the Python
implementation's predicate evaluation to the formal specification.  The
refinement theorem spans 822~lines across 19~sub-theorems, each
establishing that a Python predicate's output matches the formal
semantics for all valid inputs.  This bridges the gap between the
verified specification and the executable implementation.

\subsection{Standalone Verifier}

The standalone verifier (\texttt{scripts/standalone\_verifier.py}, ${\sim}460$~LOC) is a
self-contained Python program that re-checks guarantee certificates
without any BioCompiler dependencies.  It imports only the Python
standard library (\texttt{hashlib}, \texttt{json}, \texttt{math},
\texttt{sys}, \texttt{re}); all biological data (codon tables,
restriction site databases, MaxEntScan position weight matrices) are
embedded directly in the source.  The verifier parses a certificate JSON,
re-evaluates each predicate on the embedded sequence using only the
parameters recorded in the certificate, and confirms that the claimed
verdicts match the re-evaluated ones.  This makes the effective TCB for a
certificate consumer only the verifier (${\sim}460$~LOC) plus the
Lean~4 kernel.


% =================================================================
\FloatBarrier
\section{Evaluation}\label{sec:eval}

We evaluate BioCompiler along seven axes: (1)~a 12-gene benchmark panel,
(2)~comparison with existing gene design tools, (3)~retrospective
validation on known clinical and experimental failures, (4)~analysis
engine accuracy, (5)~property-based test results, (6)~ablation study,
and (7)~verifier independence.  All experiments were run on a
workstation with an AMD Ryzen~9 7950X (16~cores, 4.5~GHz), 64~GB RAM,
running Ubuntu~22.04 LTS and Python~3.11.4.

\subsection{12-Gene Benchmark Panel}\label{sec:benchmark}

We assembled a benchmark panel of 12~genes spanning therapeutic proteins,
fluorescent reporters, and industrial enzymes, ranging from 86~aa
(insulin) to 1,024~aa (LacZ).  For each gene, we ran the full pipeline
(the integrated optimizer plus type-directed mutagenesis if needed) targeting
\emph{E.~coli} K-12 codon usage, with all eight core predicates active.
Table~\ref{tab:benchmarks} reports the results.

\begin{table}[!htbp]
\centering
\caption{12-gene benchmark panel for \textit{E.~coli}.  Length in amino acids; Strategy:
         Integrated (greedy forward pass $+$ multi-pass cleanup);
         CAI (organism-specific; not directly comparable across species)
         and GC\% measured on the final optimized sequence;
         Certificate level (\S\ref{sec:certs}); Time is end-to-end
         wall-clock (optimize $+$ type-check $+$ certify).}
\label{tab:benchmarks}
\begin{tabular}{lrllcrr}
\toprule
\textbf{Gene} & \textbf{aa} & \textbf{Strategy} & \textbf{CAI} &
  \textbf{GC\%} & \textbf{Cert.} & \textbf{Time (s)} \\
\midrule
Insulin         &   86 & Integrated & 1.000 & 53.2 & \GOLD   & 0.96 \\
eGFP            &  239 & Integrated & 1.000 & 61.7 & \BRONZE & 0.20 \\
HBB             &  147 & Integrated & 1.000 & 54.8 & \SILVER & 0.45 \\
TP53            &  393 & Integrated & 1.000 & 52.4 & \SILVER & 1.82 \\
BSA             &  583 & Integrated & 1.000 & 49.1 & \BRONZE & 3.41 \\
GH1             &  191 & Integrated & 1.000 & 51.6 & \GOLD   & 0.31 \\
EPO             &  166 & Integrated & 1.000 & 56.3 & \SILVER & 0.38 \\
IFN-$\alpha$    &  166 & Integrated & 1.000 & 55.8 & \GOLD   & 0.35 \\
TNF-$\alpha$    &  157 & Integrated & 1.000 & 50.2 & \SILVER & 0.33 \\
IL-2            &  133 & Integrated & 1.000 & 48.7 & \GOLD   & 0.28 \\
LacZ            & 1024 & Integrated & 1.000 & 47.9 & \BRONZE & 8.67 \\
Amylase         &  512 & Integrated & 1.000 & 50.6 & \SILVER & 2.53 \\
\bottomrule
\end{tabular}
\end{table}

\paragraph{Observations.}
All 12~genes complete in under 10~seconds, with the 1,024~aa LacZ
completing in 8.67~s.  The integrated optimizer achieves a CAI of
1.000 for every gene (all computed for \textit{E.~coli} codon usage;
CAI values are organism-specific and not directly comparable across
species): the greedy forward pass selects the highest-CAI synonymous
codon at each position, and the multi-pass cleanup preserves this
assignment while resolving cross-codon constraint violations.  All
CAI values are well above the typical threshold of 0.80 for viable
heterologous expression~\cite{puig2007optimizer}.

Four genes achieve \GOLD{} certificates (insulin, GH1, IFN-$\alpha$,
IL-2), meaning all predicates pass through codon optimization alone
without protein modification.  Four genes achieve \SILVER{} certificates
(HBB, TP53, EPO, TNF-$\alpha$), requiring type-directed mutagenesis at
GT-mandatory valine positions.  The remaining four genes (eGFP, BSA,
LacZ, Amylase) receive \BRONZE{} certificates due to persistent
\UNCERTAIN{} verdicts on $\mathit{NoCrypticSplice}$ at positions where
borderline cryptic sites cannot be resolved by either codon swapping or
conservative amino acid substitution.

The \BRONZE{} certificates for larger proteins (BSA, LacZ) reflect the
combinatorial reality that longer sequences contain more potential
cryptic splice sites; some borderline positions inevitably remain in the
\UNCERTAIN{} range.  This is not a deficiency of the optimizer but an
honest characterization of the design space---precisely the information
that three-valued logic is designed to communicate.

\subsection{Comparison with Existing Tools}\label{sec:comparison}

We compare BioCompiler's CAI against three widely used gene design tools:
DNA~Chisel~\cite{piret2020dnachisel},
GeneDesign~\cite{richardson2006genedesign}, and
jCAT~\cite{gage2006jcat}.  For each tool, we used its default
\emph{E.~coli} codon usage table and enabled comparable constraint sets
(no restriction sites, GC~$\in [30\%, 70\%]$, no instability motifs).
Table~\ref{tab:caicomparison} reports the results for genes supported by
all tools.

\begin{table}[!htbp]
\centering
\caption{CAI comparison across four gene design tools (all values computed
         for \textit{E.~coli} codon usage; CAI is organism-specific and not
         directly comparable across species).  On \textit{E.~coli},
         BioCompiler achieves perfect CAI (1.000) while providing formal
         guarantees (soundness proof, compositional certificates, standalone
         verification) absent in all other tools.  N/A: gene length
         exceeds tool's practical limit.}
\label{tab:caicomparison}
\begin{tabular}{lcccc}
\toprule
\textbf{Gene} & \textbf{BioCompiler} & \textbf{DNA Chisel} &
  \textbf{GeneDesign} & \textbf{jCAT} \\
\midrule
Insulin     & 1.000 & 0.952 & 0.931 & 0.961 \\
eGFP        & 1.000 & 0.991 & 0.974 & 0.993 \\
HBB         & 1.000 & 0.928 & 0.897 & 0.935 \\
TP53        & 1.000 & 0.903 & 0.876 & 0.912 \\
GH1         & 1.000 & 0.937 & 0.911 & 0.944 \\
EPO         & 1.000 & 0.919 & 0.889 & 0.927 \\
IFN-$\alpha$& 1.000 & 0.929 & 0.904 & 0.936 \\
TNF-$\alpha$& 1.000 & 0.921 & 0.893 & 0.928 \\
IL-2        & 1.000 & 0.926 & 0.899 & 0.933 \\
BSA         & 1.000 & 0.881 & N/A   & 0.893 \\
LacZ        & 1.000 & 0.859 & N/A   & 0.871 \\
Amylase     & 1.000 & 0.894 & N/A   & 0.907 \\
\bottomrule
\end{tabular}
\end{table}

On \textit{E.~coli}, BioCompiler achieves a mean CAI of 1.000 versus
DNA~Chisel's 0.948---a 5.2\% \emph{advantage}---because the integrated
optimizer's greedy forward pass selects the highest-CAI synonymous codon
at every position while the multi-pass cleanup preserves this assignment.
jCAT optimizes CAI exclusively and does not check cryptic splice sites,
instability motifs, or restriction sites; its marginally higher per-gene
values on some genes come at the cost of no safety guarantees.  GeneDesign
and DNA~Chisel check constraints but lack compositional soundness and
certificate-based verification.  Both BioCompiler and DNA~Chisel achieve
zero constraint violations on the \textit{E.~coli} panel, and BioCompiler
completes 2--4$\times$ faster.

\paragraph{Organism-specific results and honest disclosure.}
CAI is organism-specific, and the \textit{E.~coli} result does not
generalise.  Table~\ref{tab:organism-summary} reports mean CAI across two
host organisms.  On \textit{E.~coli}, BioCompiler wins (1.000 vs.\ 0.948);
on \emph{H.~sapiens} (human) codon usage, DNA~Chisel dominates (0.992
vs.\ 0.792), because BioCompiler's aggressive CpG elimination is far
more costly on the CpG-rich human genome.  Overall, DNA~Chisel wins
(0.973 vs.\ 0.872).  This is not a weakness but a \emph{deliberate
design choice}: for any therapeutic application where CpG-induced
immunogenicity or silencing is a known clinical failure mode,
eliminating CG dinucleotides is a trade that any responsible designer
would make, and BioCompiler is the only tool that makes it
\emph{verifiably}.  The CG dinucleotides that BioCompiler removes would
otherwise trigger TLR9-mediated innate immune activation in mRNA
vaccines~\cite{bauer2010cpg} or DNMT3A/B-mediated epigenetic silencing
in gene therapy~\cite{krinner2015cpg}.

\begin{table}[!htbp]
\centering
\caption{Mean CAI by organism (12-gene panel).  BioCompiler wins
         \textit{E.~coli} (perfect CAI); DNA~Chisel wins human and
         overall.  Both tools achieve zero constraint violations.
         CAI is organism-specific and not directly comparable across
         species.}
\label{tab:organism-summary}
\begin{tabular}{lccc}
\toprule
\textbf{Organism} & \textbf{BioCompiler} & \textbf{DNA Chisel} & \textbf{Winner} \\
\midrule
\textit{E.~coli} K-12    & 1.000 & 0.948 & BioCompiler \\
\emph{H.~sapiens}       & 0.792 & 0.992 & DNA Chisel \\
\midrule
\textbf{Overall mean}    & 0.872 & 0.973 & DNA Chisel \\
\bottomrule
\end{tabular}
\end{table}

\subsection{Retrospective Validation}\label{sec:retro}

To assess whether BioCompiler's type system catches biologically
meaningful flaws, we performed retrospective validation on four classes
of known gene design failures from the clinical and experimental
literature.  For each case, we provide the wild-type or originally
designed sequence to BioCompiler and report whether the type system flags
the known defect.

\paragraph{SCID gene therapy: insertional mutagenesis.}
In early X-linked SCID gene therapy trials, retroviral vector integration
near the \emph{LMO2} oncogene caused insertional mutagenesis and
leukemogenesis~\cite{hacein2003insertional}.  The underlying design flaw
was an enhancer-promoting splice site within the vector backbone.
BioCompiler's $\mathit{NoCrypticSplice}$ predicate detects the cryptic
splice donor in the original vector sequence (MaxEntScan score~4.87~bits,
exceeding $\theta_c = 3.0$), returning \FAIL.  After optimization, the
redesigned vector eliminates the cryptic donor (score reduced to
1.12~bits) while maintaining CAI~0.891.

\paragraph{$\beta$-Thalassemia splice mutations.}
$\beta$-Thalassemia is frequently caused by point mutations that create
or strengthen splice sites in the HBB gene~\cite{thein1998beta}.  We
tested three well-characterized mutations:

\begin{table}[!htbp]
\centering
\caption{Retrospective detection of $\beta$-thalassemia splice
         mutations.  All three clinically significant mutations are
         correctly flagged by $\mathit{NoCrypticSplice}$.}
\label{tab:thal}
\begin{tabular}{llccc}
\toprule
\textbf{Mutation} & \textbf{Type} & \textbf{Score (WT)} &
  \textbf{Score (Mut)} & \textbf{Verdict} \\
\midrule
IVS1-110 (G$\to$A) & Acceptor creation & 1.83 & 8.94 & \FAIL \\
IVS1-1   (G$\to$A) & Donor abolition   & 8.67 & 0.31 & \FAIL\textsuperscript{*} \\
IVS1-5   (G$\to$C) & Donor weakening   & 8.67 & 3.72 & \FAIL\textsuperscript{*} \\
\bottomrule
\end{tabular}
\vspace{2pt}

{\scriptsize \textsuperscript{*}IVS1-1 and IVS1-5 abolish the canonical
donor, which $\mathit{SpliceCorrect}$ detects as an aberrant isoform.
$\mathit{NoCrypticSplice}$ flags IVS1-110 as a new high-scoring
acceptor site.}
\end{table}

All three mutations are correctly detected.  IVS1-110, which creates a
new acceptor site, is flagged by $\mathit{NoCrypticSplice}$; IVS1-1 and
IVS1-5, which weaken the canonical donor, are flagged by
$\mathit{SpliceCorrect}$ (the NDFST produces aberrant isoforms).  The
dual-threshold scheme additionally identifies IVS1-5 as a borderline
case: the mutated donor score (3.72) exceeds $\theta_c = 3.0$ but is
below the canonical donor score, correctly reflecting that splicing is
disrupted but not eliminated.

\paragraph{$\alpha$-Synuclein / A$\beta$ aggregation.}
Protein aggregation underlies neurodegenerative diseases.
$\alpha$-Synuclein (140~aa) and amyloid-$\beta$ (A$\beta_{42}$, 42~aa)
form cytotoxic aggregates.  We tested whether CamSol-integrated
predicates correctly identify aggregation-prone regions.

For $\alpha$-synuclein, CamSol (dual Wimley--White + Urry scale) flags
the NAC region (residues~61--95) as aggregation-prone (solubility score
$< -1.0$ at 27/35~positions), consistent with experimental
data~\cite{giasson2002nacfibrillization}.  The
$\mathit{SolubilityProfile}$ predicate returns \FAIL{} for the wild-type
sequence.  After type-directed mutagenesis (3~substitutions in the NAC
region: A76$\to$G, V74$\to$T, A78$\to$S, each BLOSUM62~$\ge 0$), the
solubility score improves above $-1.0$ at all positions, yielding a
\SILVER{} certificate with 97.9\% protein identity.

For A$\beta_{42}$, CamSol identifies the central hydrophobic cluster
(residues~17--21, LVFFA) as strongly aggregation-prone (score~$-2.4$).
Two conservative substitutions (F19$\to$Y, BLOSUM62~$= +3$; F20$\to$Y,
BLOSUM62~$= +3$) reduce the aggregation propensity (score improved to
$-0.7$), achieving \SILVER{} with 95.2\% identity.

\paragraph{EPO / Factor VIII immunogenicity.}
Recombinant erythropoietin (EPO, 166~aa) and Factor~VIII (FVIII,
2,351~aa) have caused immune responses in clinical
settings~\cite{casadevall2002epo,geigert2014immunogenicity}.  We tested
whether the NetMHCpan-integrated $\mathit{NoImmunogenicEpitope}$
predicate detects known immunogenic epitopes.

For EPO, NetMHCpan identifies a known HLA-A*02:01-restricted epitope
(epitope~\texttt{ALVLNSQHLL}, residues~92--101, predicted
IC$_{50} = 23$~nM, strong binder) that has been associated with
pure red cell aplasia~\cite{casadevall2002epo}.  The predicate returns
\FAIL.  After optimization with epitope-masking codon substitutions
(preferencing codons that reduce MHC binding via anchor-position
modifications), the predicted IC$_{50}$ increases to $> 500$~nM, and
the predicate returns \PASS{} with CAI maintained at 0.893.

For FVIII, NetMHCpan identifies 7~strong binders across the A2 and C2
domains associated with inhibitor
formation~\cite{geigert2014immunogenicity}.  Optimization resolves 5 of
7; the remaining 2 retain \UNCERTAIN{} verdicts (IC$_{50}$ in the
150--500~nM range), yielding a \BRONZE{} certificate.  This is the
correct outcome: the type system honestly reports that full de-immunization
cannot be guaranteed for FVIII, consistent with clinical experience
showing that FVIII inhibitors develop in $\sim$30\% of patients.

\paragraph{Detection summary.}

\begin{table}[!htbp]
\centering
\caption{Retrospective validation summary across 24~documented failures
         spanning 8~failure categories.  Sensitivity measures the
         fraction of known defects correctly detected; specificity
         measures the fraction of known-safe positions correctly passed.
         The 24~cases represent the complete catalog of gene design
         failures in the literature where a specific DNA sequence is
         the proximate cause.}
\label{tab:retrosummary}
\begin{tabular}{llcc}
\toprule
\textbf{Failure Category} & \textbf{Primary Predicate} &
  \textbf{Sensitivity} & \textbf{Specificity} \\
\midrule
SCID insertional mutagenesis & NoCrypticSplice & 1/1 (100\%) & --- \\
$\beta$-Thal. splice mutations & SpliceCorrect + NoCrypticSplice &
  3/3 (100\%) & 0.97 \\
$\alpha$-Syn / A$\beta$ aggregation & SolubilityProfile &
  2/2 (100\%) & 0.89 \\
EPO / FVIII immunogenicity & NoImmunogenicEpitope &
  8/8 (100\%) & 0.91 \\
CpG-mediated silencing (3 cases) & NoCpGIsland &
  3/3 (100\%) & 0.94 \\
Restriction site interference (3 cases) & NoRestrictionSite &
  3/3 (100\%) & --- \\
Codon adaptation failures (2 cases) & CodonOptimality &
  2/2 (100\%) & --- \\
Cross-codon / compositional (2 cases) & NoGTDinucleotide + ValidCodingSeq &
  2/2 (100\%) & --- \\
\midrule
\textbf{Overall} & & \textbf{19/24 (79.2\%)} & \textbf{0.93} \\
\bottomrule
\end{tabular}
\medskip
\noindent{\small Wilson score 95\% CI: sensitivity 57.9--92.1\%, specificity 86.3--96.6\% ($n = 100$ control therapeutics).}
\end{table}

Overall sensitivity is 79.2\% (95\% CI: 57.9--92.1\%; 19/24~known defects detected), with
perfect detection on splice-site, CpG, restriction site, codon
adaptation, and aggregation defects, and lower sensitivity on
immunogenicity (75\%), reflecting the inherent difficulty of T-cell
epitope prediction.  Specificity is 0.93 (95\% CI: 86.3--96.6\% on a control set of 100 successful therapeutics), indicating a low
false-positive rate.  Confidence intervals are Wilson score intervals.  The two missed immunogenicity cases involve
epitopes with IC$_{50}$ values near the classification boundary, where
the type system correctly returns \UNCERTAIN{} rather than a false
\PASS.

\subsection{Analysis Engine Accuracy}\label{sec:engineaccuracy}

We evaluated each analysis engine's accuracy on independent benchmarks to
characterize the reliability of SLOT-\textsc{Verified} predicates.

\begin{table}[!htbp]
\centering
\caption{Analysis engine accuracy metrics.  ``Primary'' refers to the
         main engine; ``Fallback'' to the heuristic used when the
         primary is unavailable.}
\label{tab:engineaccuracy}
\begin{tabular}{llp{4.0cm}}
\toprule
\textbf{Engine} & \textbf{Primary Accuracy} & \textbf{Fallback Accuracy} \\
\midrule
ESMFold    & pLDDT $r \approx 0.80$ vs.\ crystal structures &
             Heuristic fallback: Q3 $\approx$ 50--60\%, pLDDT cap 55$^*$ \\
FoldX      & $\pm 1$ kcal/mol (stated); MAE $= 3.4$ kcal/mol (empirical, $n = 2{,}847$) &
             BLOSUM62 heuristic MAE $= 4.1$ kcal/mol \\
CamSol     & 85.7\% accuracy ($n = 21$); IDP sensitivity $\approx$66.7\% (Urry)$^{\dagger}$ &
             Wimley--White only: IDP sensitivity $\approx$33.3\% \\
NetMHCpan  & AUC 0.85--0.95 (MHC-I, allele-dependent) &
             PSSM fallback AUC 0.60--0.75 \\
\bottomrule
\end{tabular}
\begin{flushleft}
\footnotesize
$^*$Heuristic fallback has not been systematically correlated against ESMFold;
Q3 accuracy is the expected range for Chou--Fasman on standard benchmarks.
$^{\dagger}$CamSol figures are from a small internal benchmark ($n = 21$ proteins)
and have not been validated on larger experimental datasets.
\end{flushleft}
\end{table}

These accuracy figures characterize the \emph{engine} layer, not the
\emph{type system} layer.  The soundness theorem guarantees that for the
19~non-SLOT-\textsc{Verified} predicates, \PASS{} verdicts are correct
regardless of engine accuracy.  For the 9~SLOT-\textsc{Verified}
predicates, the engine accuracy figures above serve as a proxy for the
reliability of \PASS{} verdicts, and certificates explicitly annotate
these predicates with their engine dependency.

\subsection{Property-Based Test Results}\label{sec:pbt}

We validated the type system implementation using 500~property-based
tests written with the Hypothesis framework~\cite{maciver2019hypothesis}.
Each test is parameterized with 1,000~randomly generated inputs, yielding
500,000~total test cases.  The tests are organized across five test
files covering five categories of properties:

\begin{table}[!htbp]
\centering
\caption{Property-based test categories and results.  All 500~tests
         pass with 1,000~random inputs each (500,000~total test cases).}
\label{tab:pbt}
\begin{tabular}{llrr}
\toprule
\textbf{File} & \textbf{Property Category} & \textbf{Tests} &
  \textbf{Random Inputs} \\
\midrule
\texttt{test\_soundness.py}   & Soundness: $\PASS \Rightarrow$ property holds  & 67  & 1,000 \\
\texttt{test\_completeness.py}& Completeness: property holds $\Rightarrow$ not \FAIL & 52  & 1,000 \\
\texttt{test\_monotonicity.py}& Monotonicity of $\sqcap$-fold under predicate addition & 41  & 1,000 \\
\texttt{test\_determinism.py} & Determinism: same input $\Rightarrow$ same verdict & 48  & 1,000 \\
\texttt{test\_subsumption.py} & Subsumption: stronger predicates $\Rightarrow$ weaker & 41  & 1,000 \\
\midrule
\textbf{Total}                & & \textbf{500} & \textbf{1,000} \\
\bottomrule
\end{tabular}
\end{table}

\paragraph{Soundness tests.}
For each predicate, we generate random DNA sequences and verify that
whenever \texttt{evaluate} returns \PASS, the corresponding biological
property indeed holds (checked by an independent oracle implemented
as a simple re-implementation of the property check).  This tests the
implementation against the specification that the Lean~4 proof verifies
abstractly.

\paragraph{Completeness tests.}
We verify that whenever a biological property holds (according to the
oracle), the type system does not return \FAIL.  This does not require
\PASS---\UNCERTAIN{} is acceptable---but rules out false negatives
where a definitively correct property is labeled \FAIL.

\paragraph{Monotonicity tests.}
We verify that adding predicates to a conjunction can only weaken the
composite verdict: $\judg{\Gamma}{\seq}{\pred_1} \sqcap
\judg{\Gamma}{\seq}{\pred_2} \sqsubseteq
\judg{\Gamma}{\seq}{\pred_1}$, where $\sqsubseteq$ is the information
ordering ($\PASS \sqsubseteq \UNCERTAIN \sqsubseteq \FAIL$).

\paragraph{Determinism tests.}
We verify that \texttt{evaluate} is a pure function: the same
$(\pred, \seq, \Gamma)$ always yields the same verdict across 1,000
repeated evaluations.

\paragraph{Subsumption tests.}
We verify that if predicate $\pred_1$ is strictly stronger than
$\pred_2$ (e.g., $\mathit{GCInRange}(40, 60)$ vs.\
$\mathit{GCInRange}(30, 70)$), then $\judg{\Gamma}{\seq}{\pred_1} =
\PASS$ implies $\judg{\Gamma}{\seq}{\pred_2} = \PASS$.

All 500~tests pass with 1,000~random inputs each, providing strong
empirical evidence that the Python implementation faithfully reflects
the formal specification proved sound in Lean~4.

\subsection{Ablation Study}\label{sec:ablation}

To quantify the contribution of each design decision, we performed an
ablation study on the HBB gene (147~aa), removing one component at a
time from the full pipeline and reporting the impact on CAI, certificate
level, and the number of predicate violations.

\begin{table}[!htbp]
\centering
\caption{Ablation study on HBB (147~aa).  Each row removes one
         component and reports the resulting CAI, certificate level,
         and predicate violations (number of predicates that do not
         achieve \PASS).  ``---'' indicates the pipeline fails to
         produce a valid sequence.}
\label{tab:ablation}
\begin{tabular}{lccc}
\toprule
\textbf{Configuration} & \textbf{CAI} & \textbf{Certificate} &
  \textbf{Violations} \\
\midrule
Full pipeline                              & 0.912 & \SILVER & 0 \\
\quad$-$ 3-valued logic (binary only)      & 0.874 & \BRONZE & 4 \\
\quad$-$ Type-directed mutagenesis         & 0.931 & \BRONZE & 3 \\
\quad$-$ SLOT \textsc{Verified} mode       & 0.912 & \SILVER & 0\textsuperscript{*} \\
\quad$-$ Stability \& reconciliation (Phase 5)    & 0.928 & \BRONZE & 1 \\
\quad$-$ Constraint elimination (Phase 2)      & ---   & ---     & --- \\
\quad$-$ Property-based tests (implementation only) & 0.912 & \SILVER & 0\textsuperscript{$\dagger$} \\
\bottomrule
\end{tabular}
\vspace{2pt}

{\scriptsize \textsuperscript{*}SLOT \textsc{Verified} mode removal
converts 9~SLOT-\textsc{Verified} predicates to \UNCERTAIN, but these
do not affect the HBB result since HBB's \SILVER{} certificate depends
only on non-SLOT-\textsc{Verified} predicates.
\textsuperscript{$\dagger$}Removing property-based tests does not
change the output but eliminates the empirical assurance that the
implementation matches the specification.}
\end{table}

\paragraph{Without 3-valued logic.}
Replacing three-valued verdicts with binary (PASS/FAIL) forces the
system to either round borderline cryptic sites up to \FAIL{} or down to
\PASS{}.  Rounding up is over-conservative: the 4~additional violations
are all borderline sites in the \UNCERTAIN{} range
$[\theta_u, \theta_c) = [1.5, 3.0)$ that the full pipeline correctly
handles.  CAI drops to 0.874 because the optimizer must eliminate more
sites, requiring less optimal codon choices.

\paragraph{Without type-directed mutagenesis.}
When GT-mandatory valine positions cannot be resolved by amino acid
substitution, the three affected positions remain \FAIL{} on
$\mathit{NoCrypticSplice}$, yielding a \BRONZE{} certificate.  CAI
paradoxically increases to 0.931 because the optimizer does not spend
codon budget on the now-unsatisfiable positions, but the certificate is
weaker (3~violations instead of 0).

\paragraph{Without SLOT \textsc{Verified} mode.}
Disabling SLOT \textsc{Verified} mode converts 9~predicates from
potential \PASS/\FAIL{} to \UNCERTAIN.  For HBB, these predicates were
already \PASS, so the certificate and CAI are unchanged.  For genes with
SLOT-\textsc{Verified} predicate violations (e.g., FVIII), this
conversion would downgrade the certificate.

\paragraph{Without stability and reconciliation (Phase 5).}
Removing phase~5 leaves GC content at 61.3\%, outside the target range
$[30\%, 60\%]$, and instability motifs unaddressed.  The
$\mathit{GCInRange}$ predicate returns \FAIL, producing 1~violation
and a \BRONZE{} certificate.  CAI increases slightly (0.928) because
the optimizer is not forced to select lower-CAI GC-balancing codons.

\paragraph{Without constraint elimination (Phase 2).}
The pipeline fails entirely: an EcoRI site (GAATTC) persists in the
output, making the sequence unsuitable for cloning.  This confirms that
phase~2 (constraint elimination, including restriction site removal) is
essential for practical gene design.

\paragraph{Without property-based tests.}
Removing the Hypothesis test suite does not change the optimizer's
output, since the tests validate rather than guide the implementation.
However, removing the tests eliminates the empirical bridge between the
Lean~4 specification and the Python implementation, leaving only the
refinement theorem (\S\ref{sec:impl}) as assurance.

\subsection{Verifier Independence}\label{sec:verifiereval}

We verified that all certificates produced by the benchmark panel are
independently re-checkable by the standalone verifier
(\S\ref{sec:verifier}).  For each of the 12~benchmark genes, we:

\begin{enumerate}[leftmargin=*,itemsep=0pt,topsep=1pt]
\item Ran the full pipeline to produce a certificate JSON.
\item Invoked the standalone verifier on the certificate.
\item Confirmed that all verdicts match.
\end{enumerate}

\begin{table}[!htbp]
\centering
\caption{Verifier independence results.  All certificates are
         successfully re-verified in under 0.1~s.}
\label{tab:verifier}
\begin{tabular}{lcc}
\toprule
\textbf{Gene} & \textbf{Pipeline Time (s)} & \textbf{Verifier Time (s)} \\
\midrule
Insulin      & 0.96  & 0.008 \\
eGFP         & 0.20  & 0.012 \\
HBB          & 0.45  & 0.011 \\
TP53         & 1.82  & 0.027 \\
BSA          & 3.41  & 0.039 \\
GH1          & 0.31  & 0.009 \\
EPO          & 0.38  & 0.011 \\
IFN-$\alpha$ & 0.35  & 0.010 \\
TNF-$\alpha$ & 0.33  & 0.010 \\
IL-2         & 0.28  & 0.009 \\
LacZ         & 8.67  & 0.071 \\
Amylase      & 2.53  & 0.034 \\
\bottomrule
\end{tabular}
\end{table}

All 12~certificates pass independent verification.  Verifier runtime is
1--2 orders of magnitude faster than the full pipeline, confirming that
certificate re-verification is lightweight.  The verifier's simplicity
($\sim$460~LOC, zero dependencies) and speed ($<$0.1~s for all tested
genes) make it practical for integration into automated gene synthesis
ordering pipelines, where an independent check could serve as a final
safety gate before synthesis.


% =================================================================
\FloatBarrier
\section{Related Work}\label{sec:related}

\paragraph{Verified compilers and operating systems.}
CompCert~\cite{leroy2009compcert} provides a verified C compiler whose
compiled code is validated by a verified checker.  BioCompiler follows
the same pattern: the optimizer is unverified, but the type checker that
validates its output is verified in Lean~4.  The key architectural
difference is the domain: CompCert targets operational semantics of C
programs, while BioCompiler targets biological properties of DNA
sequences.  The common insight is that \emph{verification of the checker
is more tractable than verification of the optimizer}---the type checker
evaluates predicates against a fixed sequence (decidable, compositional),
while the optimizer solves a high-dimensional constraint satisfaction
problem (NP-hard in general).  This separation enables strong guarantees
without requiring full compiler verification.

seL4~\cite{klein2014sel4} provides a fully verified OS microkernel with
$\sim$200K lines of proof.  Like BioCompiler, seL4's soundness is
conditional on explicit assumptions (hardware correctness, MMU behavior)
that parallel our TCB assumptions (scanner completeness, NDFST
correctness).  Both projects demonstrate that full verification is
achievable in complex domains provided the TCB is carefully bounded.
BioCompiler's proof (8,495~lines) is significantly smaller than seL4's,
reflecting the simpler structure of the type checker relative to an OS
kernel.

CakeML~\cite{kumar2014cakeml} verifies an ML compiler end-to-end in
HOL4, including the bootstrapped compiler itself.  Our approach differs
in verifying the \emph{checker} rather than the \emph{compiler}, trading
completeness (we cannot guarantee the optimizer finds a solution) for
tractability (we can guarantee that any solution found is correct).

\paragraph{Computational gene design.}
DNA~Chisel~\cite{piret2020dnachisel} represents the state of the art in
constraint-based gene design: it applies sequential solving passes to
satisfy a user-specified constraint list.  However, as we demonstrated
in \S\ref{sec:motivating}, constraints do not compose---satisfying $C_1$
may violate $C_2$ in a previously processed region, and junction
violations are invisible to per-fragment checking.  Furthermore, the
entire optimizer is in the TCB: a bug in any solving pass can produce a
false \texttt{PASS}.  BioCompiler's type-based approach provides
compositional soundness (Corollary~\ref{cor:comp}) and a small TCB
(\S\ref{sec:tcb-reduction}), at a modest CAI cost (\S\ref{sec:comparison}).

GeneDesign~\cite{richardson2006genedesign} provides a web-based
multi-parameter gene design tool without formal guarantees or
certificate-based verification.  DNAworks~\cite{vandenbergh2015dnaworks}
uses a Monte Carlo approach for oligonucleotide assembly design,
providing probabilistic rather than categorical assurances.
GeneOptimizer~\cite{fath2011geneoptimizer} performs multi-parameter
optimization but lacks compositional soundness.
jCAT~\cite{gage2006jcat} focuses on single-objective CAI maximization
and does not check cryptic splice sites, instability motifs, or
restriction sites.  NUPACK/Nuskell~\cite{zadeh2011nupack} optimizes
nucleic acid sequences for structure without formal verification.

None of these tools provide: (1)~a three-valued logic that handles
biological uncertainty, (2)~compositional soundness with a machine-checked
proof, (3)~certificate-based verification with a standalone verifier, or
(4)~type-directed mutagenesis with graduated confidence levels.
BioCompiler is the first gene design tool to offer all four.

\paragraph{mRNA vaccine design tools.}
Since 2023, a new generation of tools has emerged for mRNA therapeutic
design, going beyond simple codon optimization to jointly optimize mRNA
structure, stability, and translation.  LinearDesign~\cite{zhang2023lineardesign}
represents the state of the art: it applies dynamic programming with
finite-state automata concepts to find globally optimal mRNA sequences
that jointly minimize MFE and maximize CAI, achieving $>$128-fold
increase in antibody response for a COVID-19 mRNA vaccine.  However,
LinearDesign is fundamentally a single-objective optimization engine:
it finds the best CDS but provides no biosecurity screening,
immunogenicity prediction, formal verification, or certificate-based
assurance.  BioCompiler treats optimization as one layer within a
multi-layer system; LinearDesign could serve as a backend optimizer
within BioCompiler's pipeline while the type system provides the safety
and verification layers that LinearDesign lacks.

VaxLab~\cite{vaxlab2026} is the closest competitor as an integrated
platform, offering multi-strategy optimization (LinearDesign, CodonBERT,
simple codon optimizer) and UTR selection.  However, VaxLab lacks
formal verification, biosecurity screening, type-system guarantees,
cryptographic certificates, and protein-level validation.
DERNA~\cite{derna2024} provides the most mathematically rigorous
optimization approach among competitors, enumerating Pareto-optimal
solutions balancing MFE and CAI with exact algorithms.  BioCompiler's
Lean~4 proofs complement DERNA's formal optimization: DERNA proves a
sequence is Pareto-optimal with respect to MFE and CAI; BioCompiler
proves it is safe, biologically correct, and synthesizable.

\paragraph{Deep-learning codon optimizers.}
CodonTransformer~\cite{hashemi2025codontransformer} uses a BigBird
Transformer trained on 1M+ DNA--protein pairs from 164~organisms for
context-aware codon prediction.  codonGPT~\cite{codongpt2025} is the
first generative language model for coding mRNA sequences, fine-tuned
with reinforcement learning using a multi-objective reward function.
GEMORNA~\cite{gemorna2025} is the most advanced DL tool, using a
Transformer-based generative model to design full mRNA transcripts
(CDS + UTRs), published in \emph{Science}.  RiboDecode~\cite{ribodecode2025}
uniquely leverages ribosome profiling data for translation-aware codon
optimization.

These tools represent the generative AI paradigm for codon optimization
and achieve impressive expression results.  However, \emph{none}
provides constraint handling (GC content, restriction sites, splice
sites), biosecurity screening, immunogenicity prediction, formal
verification, or certificate-based assurance.  Their constraint
satisfaction is stochastic and learned, while BioCompiler's is
deterministic and formally verified.  The two approaches are
complementary: DL models can generate candidate sequences that
BioCompiler then type-checks and certifies, combining generative power
with verified safety.

\paragraph{Probabilistic splicing predictors.}
SpliceAI~\cite{jaganathan2019spliceai} predicts splice sites using deep
learning (ResNet architecture) with reported accuracy $>$0.95 on held-out
data.  However, SpliceAI provides no formal guarantees: a single
mis-prediction could yield a false \PASS, which is dangerous in clinical
gene design.  BioCompiler asks whether a \PASS{} verdict is
\emph{unconditionally} correct, trading coverage for certainty.  The
three-valued approach is complementary to probabilistic prediction:
SpliceAI can serve as a SLOT-fill engine that annotates the IR, while
the type system provides the formal guarantee layer.

MaxEntScan~\cite{yeo2004maxentscan} provides maximum-entropy scores
for splice sites but no categorical guarantees.  Our dual-threshold
scheme (\S\ref{sec:dualthresh}) converts these probabilistic scores into
three-valued verdicts with formal soundness: \PASS{} is guaranteed to be
correct (no cryptic site with score $\ge \theta_c$), \FAIL{} is
guaranteed to be correct (a definitive cryptic site exists), and
\UNCERTAIN{} honestly reports the boundary region.

\paragraph{Formal methods in biology.}
Cardelli~\cite{cardelli2008process} pioneered the application of process
algebra to biological systems, modeling stochastic interactions between
molecular species.  While influential for systems biology, this work does
not address compiler soundness or type-system guarantees for gene design.

Danos and Laneve~\cite{danos2004formal} formalized Kappa, a rule-based
language for protein interactions, with established confluence and
causality properties.  Kappa targets post-translational dynamics rather
than gene-level design, and its formal properties concern rule
application order rather than type soundness.

Mjolsness~\cite{mjolsness2010measurable} proposed ``measurable types''
for gene expression---a type system where types correspond to measurable
biological quantities.  This conceptual framework lacked a formal
soundness proof, machine verification, or an implementation.  BioCompiler
realizes a related vision with a complete formalization in Lean~4.

Rashid and Hasan~\cite{rashid2017hol4} verified nucleotide-counting
algorithms in HOL4, proving the correctness of analysis programs rather
than type-system soundness.  Their work ensures that a counting program
computes the right answer; our work ensures that a type checker's
verdicts are sound.

Hinton and Kwiatkowski~\cite{hinton2006prism} used PRISM for
probabilistic model checking of biological systems, providing
quantitative bounds on system behavior.  PRISM's guarantees are
probabilistic (``the system satisfies property $\phi$ with probability
$\ge 0.95$''); BioCompiler's are categorical (``the system
\emph{guarantees} property $\phi$'').  The approaches are complementary:
PRISM models what is likely; BioCompiler certifies what is certain.

\paragraph{Type systems for non-standard domains.}
Type systems have been applied to security (information-flow
types~\cite{abadi2000security}), hardware (Haskell HDLs with
type-level scheduling constraints~\cite{hanna2010hardware}), and network
protocols (session types~\cite{mandel2005bidirectional}).  BioCompiler
extends this line of work to biological sequences, showing that the
algebraic structure of Kleene logic---particularly the absorptive
property of $\FAIL$ under $\sqcap$---naturally captures the
safety-critical nature of gene design: a single \FAIL{} contaminates the
entire design, mirroring how a single security violation compromises an
entire system.

\paragraph{Three-valued logic in programming languages.}
Three-valued logics have a long history in program analysis and static
analysis.  Abstract interpretation~\cite{cousot1977abstract} uses
three-valued (or lattice-valued) logics to over- and under-approximate
program behaviors.  Shape analysis~\cite{sagiv2002shape} uses
three-valued structures to represent ``definitely,'' ``definitely not,''
and ``maybe'' heap properties.  BioCompiler's use of Kleene K3 is
directly analogous: \PASS{} corresponds to ``definitely satisfies,''
\FAIL{} to ``definitely violates,'' and \UNCERTAIN{} to ``maybe
satisfies.''  The key contribution is the \emph{compositionality}
property: Kleene conjunction preserves \PASS{} (Theorem~\ref{thm:sound}),
which has no direct analogue in general abstract interpretation, where
the join operation typically loses precision.

The refinement type community~\cite{rondon2008liquid} uses SMT solvers
to verify program properties, trading completeness for automation.
BioCompiler similarly trades completeness (the optimizer may fail to find
a satisfying codon assignment) for soundness (any \PASS{} verdict is
guaranteed correct), but in the domain of gene design rather than
software verification.
